\documentclass{article}
\usepackage[utf8]{inputenc}
\usepackage[T1]{fontenc}
\usepackage[paperwidth=320mm,paperheight=180mm,margin=0mm]{geometry}
\usepackage{tikz,amsmath,amssymb,tgheros,hyperref}
\renewcommand{\familydefault}{\sfdefault}
\pagestyle{empty}
\setlength{\parindent}{0pt}
\definecolor{Navy}{HTML}{102A43}
\definecolor{Blue}{HTML}{2563EB}
\definecolor{Gray}{HTML}{52606D}
\definecolor{Pale}{HTML}{F3F6FA}
\definecolor{Ice}{HTML}{DBEAFE}
\definecolor{Line}{HTML}{D7E2EB}
\hypersetup{pdfauthor={Jason Chambe},pdftitle={Tamaño de muestra en R: nueve variantes},colorlinks=true,urlcolor=Blue}
\begin{document}
\noindent\begin{tikzpicture}[remember picture,overlay,x=1mm,y=1mm]\begin{scope}[shift={(current page.north west)}]
\fill[Navy,rounded corners=0mm] (0,-0) rectangle (320,-180);
\node[anchor=north west,inner sep=0pt,text width=265mm,align=left,text=Ice,font=\fontsize{10}{12.2}\selectfont\bfseries] at (14,-10) {TUTORIAL REPRODUCIBLE EN R};
\node[anchor=north west,inner sep=0pt,text width=24mm,align=left,text=white,font=\fontsize{10}{12.2}\selectfont\bfseries] at (286,-10) {01 / 20};
\draw[Gray,line width=.4pt] (14,-169) -- (306,-169);
\node[anchor=north west,inner sep=0pt,text width=250mm,align=left,text=Ice,font=\fontsize{9}{11.0}\selectfont] at (14,-172) {Jason Chambe  ·  Bioestadística  ·  TAREA 2};
\node[anchor=north west,inner sep=0pt,text width=34mm,align=left,text=Ice,font=\fontsize{9}{11.0}\selectfont\bfseries] at (274,-172) {DESLIZA >};
\node[anchor=north west,inner sep=0pt,text width=198mm,align=left,text=white,font=\fontsize{43}{47}\selectfont\bfseries] at (14,-37) {Tamaño de\\muestra en R};
\node[anchor=north west,inner sep=0pt,text width=174mm,align=left,text=Ice,font=\fontsize{23}{28.1}\selectfont] at (15,-88) {Entender qué comparamos antes de elegir una función.};
\fill[Blue,rounded corners=3mm] (213,-40) rectangle (306,-144);
\node[anchor=north west,inner sep=0pt,text width=74mm,align=left,text=white,font=\fontsize{55}{67.1}\selectfont\bfseries] at (222,-48) {3};
\node[anchor=north west,inner sep=0pt,text width=71mm,align=left,text=white,font=\fontsize{16}{19.5}\selectfont\bfseries] at (223,-73) {ejercicios centrales};
\node[anchor=north west,inner sep=0pt,text width=74mm,align=left,text=white,font=\fontsize{55}{67.1}\selectfont\bfseries] at (222,-94) {9};
\node[anchor=north west,inner sep=0pt,text width=73mm,align=left,text=white,font=\fontsize{16}{19.5}\selectfont\bfseries] at (223,-119) {variantes explicadas};
\node[anchor=north west,inner sep=0pt,text width=184mm,align=left,text=white,font=\fontsize{14}{17.1}\selectfont\bfseries] at (15,-125) {Hemoglobina · Anemia · Dos distritos};
\node[anchor=north west,inner sep=0pt,text width=181mm,align=left,text=Ice,font=\fontsize{12}{14.6}\selectfont] at (15,-140) {Código, supuestos y resultados.\\Práctica completa: 8 escenarios y 24 variantes.};
\end{scope}\end{tikzpicture}\mbox{}\newpage
\noindent\begin{tikzpicture}[remember picture,overlay,x=1mm,y=1mm]\begin{scope}[shift={(current page.north west)}]
\fill[Pale,rounded corners=0mm] (0,-0) rectangle (320,-180);
\node[anchor=north west,inner sep=0pt,text width=265mm,align=left,text=Blue,font=\fontsize{10}{12.2}\selectfont\bfseries] at (14,-10) {PLANTEAR  /  PUNO · VARIANTE A · BILATERAL};
\node[anchor=north west,inner sep=0pt,text width=24mm,align=left,text=Navy,font=\fontsize{10}{12.2}\selectfont\bfseries] at (286,-10) {02 / 20};
\draw[Line,line width=.4pt] (14,-169) -- (306,-169);
\node[anchor=north west,inner sep=0pt,text width=250mm,align=left,text=Gray,font=\fontsize{9}{11.0}\selectfont] at (14,-172) {Jason Chambe  ·  Bioestadística  ·  TAREA 2};
\node[anchor=north west,inner sep=0pt,text width=34mm,align=left,text=Gray,font=\fontsize{9}{11.0}\selectfont\bfseries] at (274,-172) {DESLIZA >};
\node[anchor=north west,inner sep=0pt,text width=290mm,align=left,text=Navy,font=\fontsize{27}{32.9}\selectfont\bfseries] at (14,-25) {Hemoglobina: una media frente a referencia};
\fill[white,rounded corners=3mm] (14,-55) rectangle (191,-147);
\fill[white,rounded corners=3mm] (197,-55) rectangle (306,-147);
\node[anchor=north west,inner sep=0pt,text width=160mm,align=left,text=Blue,font=\fontsize{10}{12.2}\selectfont\bfseries] at (21,-62) {LA COMPARACIÓN};
\node[anchor=north west,inner sep=0pt,text width=160mm,align=left,text=Navy,font=\fontsize{18}{23}\selectfont] at (21,-73) {Después de la intervención esperamos una media de 14 g/dL. ¿Cuántos niños necesitamos para compararla con 10 g/dL?};
\node[anchor=north west,inner sep=0pt,text width=160mm,align=left,text=Blue,font=\fontsize{25}{30.5}\selectfont\bfseries] at (21,-108) {$10\ \mathrm{g/dL}\ \longrightarrow\ 14\ \mathrm{g/dL}$};
\node[anchor=north west,inner sep=0pt,text width=160mm,align=left,text=Navy,font=\fontsize{12.4}{15.1}\selectfont] at (21,-129) {$H_0:\mu=10\qquad H_A:\mu\ne10$};
\node[anchor=north west,inner sep=0pt,text width=95mm,align=left,text=Blue,font=\fontsize{10}{12.2}\selectfont\bfseries] at (204,-62) {SUPUESTOS DEL CÁLCULO};
\node[anchor=north west,inner sep=0pt,text width=95mm,align=left,text=Navy,font=\fontsize{14}{17.1}\selectfont] at (204,-76.0) {Referencia fija: 10 g/dL};
\node[anchor=north west,inner sep=0pt,text width=95mm,align=left,text=Navy,font=\fontsize{14}{17.1}\selectfont] at (204,-90.5) {DE esperada: 1 g/dL};
\node[anchor=north west,inner sep=0pt,text width=95mm,align=left,text=Navy,font=\fontsize{14}{17.1}\selectfont\bfseries] at (204,-105.0) {Alfa: 0,05 · Potencia: 80 \%};
\node[anchor=north west,inner sep=0pt,text width=95mm,align=left,text=Navy,font=\fontsize{14}{17.1}\selectfont] at (204,-119.5) {Alternativa: bilateral};
\fill[Ice,rounded corners=3mm] (14,-151) rectangle (306,-165);
\node[anchor=north west,inner sep=0pt,text width=282mm,align=left,text=Navy,font=\fontsize{12}{14}\selectfont] at (19,-155) {Hay una muestra nueva. La referencia de 10 se considera fija, sin incertidumbre muestral.};
\node[anchor=north west,inner sep=0pt,text width=292mm,align=left,text=Gray,font=\fontsize{8.5}{10.4}\selectfont] at (14,-45) {Teoría 01, pp. 7 y 29 · Teoría 02, pp. 9 y 11 · Teoría 03, pp. 30 y 31};
\end{scope}\end{tikzpicture}\mbox{}\newpage
\noindent\begin{tikzpicture}[remember picture,overlay,x=1mm,y=1mm]\begin{scope}[shift={(current page.north west)}]
\fill[Pale,rounded corners=0mm] (0,-0) rectangle (320,-180);
\node[anchor=north west,inner sep=0pt,text width=265mm,align=left,text=Blue,font=\fontsize{10}{12.2}\selectfont\bfseries] at (14,-10) {CALCULAR E INTERPRETAR  /  PUNO · VARIANTE A · BILATERAL};
\node[anchor=north west,inner sep=0pt,text width=24mm,align=left,text=Navy,font=\fontsize{10}{12.2}\selectfont\bfseries] at (286,-10) {03 / 20};
\draw[Line,line width=.4pt] (14,-169) -- (306,-169);
\node[anchor=north west,inner sep=0pt,text width=250mm,align=left,text=Gray,font=\fontsize{9}{11.0}\selectfont] at (14,-172) {Jason Chambe  ·  Bioestadística  ·  TAREA 2};
\node[anchor=north west,inner sep=0pt,text width=34mm,align=left,text=Gray,font=\fontsize{9}{11.0}\selectfont\bfseries] at (274,-172) {DESLIZA >};
\node[anchor=north west,inner sep=0pt,text width=292mm,align=left,text=Navy,font=\fontsize{28}{34.2}\selectfont\bfseries] at (14,-25) {Del código al tamaño que reportamos};
\node[anchor=north west,inner sep=0pt,text width=292mm,align=left,text=Blue,font=\fontsize{15}{18.3}\selectfont] at (14,-44) {$d=(14-10)/1=4$};
\fill[Navy,rounded corners=3mm] (14,-58) rectangle (194,-147);
\fill[white,rounded corners=3mm] (200,-58) rectangle (306,-147);
\node[anchor=north west,inner sep=0pt,text width=163mm,align=left,text=Ice,font=\fontsize{10}{12.2}\selectfont\bfseries] at (21,-64) {R · PAQUETE pwr};
\node[anchor=north west,inner sep=0pt,text width=166mm,align=left,text=white,font=\fontsize{13.5}{16.9}\selectfont] at (21,-76) {\ttfamily library(pwr)\\media\_referencia\ <-\ 10\\media\_esperada\ <-\ 14\\de\_esperada\ <-\ 1\\d\ <-\ (media\_esperada\ -\ media\_referencia)\ /\ de\_esperada\\calculo\ <-\ pwr.t.test(d\ =\ d,\ sig.level\ =\ 0.05,\\\ \ power\ =\ 0.80,\ type\ =\ "one.sample",\\\ \ alternative\ =\ "two.sided")\\calculo\$n\\ceiling(calculo\$n)};
\node[anchor=north west,inner sep=0pt,text width=90mm,align=left,text=Blue,font=\fontsize{10}{12.2}\selectfont\bfseries] at (208,-64) {RESULTADO};
\node[anchor=north west,inner sep=0pt,text width=91mm,align=left,text=Blue,font=\fontsize{51}{62.2}\selectfont\bfseries] at (207,-72) {3};
\node[anchor=north west,inner sep=0pt,text width=90mm,align=left,text=Navy,font=\fontsize{14}{17.1}\selectfont\bfseries] at (208,-95) {niños en una muestra};
\node[anchor=north west,inner sep=0pt,text width=90mm,align=left,text=Gray,font=\fontsize{11.5}{14.0}\selectfont] at (208,-108) {n calculado: 2,720139};
\node[anchor=north west,inner sep=0pt,text width=90mm,align=left,text=Navy,font=\fontsize{12}{14.6}\selectfont\bfseries] at (208,-120) {Redondeo hacia arriba};
\node[anchor=north west,inner sep=0pt,text width=90mm,align=left,text=Gray,font=\fontsize{11}{13.4}\selectfont] at (208,-134) {Potencia con n: 90,85 \%};
\fill[Ice,rounded corners=3mm] (14,-151) rectangle (306,-165);
\node[anchor=north west,inner sep=0pt,text width=282mm,align=left,text=Navy,font=\fontsize{12}{14}\selectfont] at (19,-154.5) {El tamaño pequeño se explica por d = 4. Depende de normalidad y de una DE bien sustentada.};
\end{scope}\end{tikzpicture}\mbox{}\newpage
\noindent\begin{tikzpicture}[remember picture,overlay,x=1mm,y=1mm]\begin{scope}[shift={(current page.north west)}]
\fill[Pale,rounded corners=0mm] (0,-0) rectangle (320,-180);
\node[anchor=north west,inner sep=0pt,text width=265mm,align=left,text=Blue,font=\fontsize{10}{12.2}\selectfont\bfseries] at (14,-10) {PLANTEAR  /  PUNO · VARIANTE B · INCREMENTO};
\node[anchor=north west,inner sep=0pt,text width=24mm,align=left,text=Navy,font=\fontsize{10}{12.2}\selectfont\bfseries] at (286,-10) {04 / 20};
\draw[Line,line width=.4pt] (14,-169) -- (306,-169);
\node[anchor=north west,inner sep=0pt,text width=250mm,align=left,text=Gray,font=\fontsize{9}{11.0}\selectfont] at (14,-172) {Jason Chambe  ·  Bioestadística  ·  TAREA 2};
\node[anchor=north west,inner sep=0pt,text width=34mm,align=left,text=Gray,font=\fontsize{9}{11.0}\selectfont\bfseries] at (274,-172) {DESLIZA >};
\node[anchor=north west,inner sep=0pt,text width=290mm,align=left,text=Navy,font=\fontsize{27}{32.9}\selectfont\bfseries] at (14,-25) {La dirección se decide antes de medir};
\fill[white,rounded corners=3mm] (14,-55) rectangle (191,-147);
\fill[white,rounded corners=3mm] (197,-55) rectangle (306,-147);
\node[anchor=north west,inner sep=0pt,text width=160mm,align=left,text=Blue,font=\fontsize{10}{12.2}\selectfont\bfseries] at (21,-62) {LA COMPARACIÓN};
\node[anchor=north west,inner sep=0pt,text width=160mm,align=left,text=Navy,font=\fontsize{18}{23}\selectfont] at (21,-73) {Con los mismos valores de hemoglobina, ahora planteamos formalmente una hipótesis unilateral de incremento.};
\node[anchor=north west,inner sep=0pt,text width=160mm,align=left,text=Blue,font=\fontsize{25}{30.5}\selectfont\bfseries] at (21,-108) {$10\ \mathrm{g/dL}\ \longrightarrow\ 14\ \mathrm{g/dL}$};
\node[anchor=north west,inner sep=0pt,text width=160mm,align=left,text=Navy,font=\fontsize{12.4}{15.1}\selectfont] at (21,-129) {$H_0:\mu\le10\qquad H_A:\mu>10$};
\node[anchor=north west,inner sep=0pt,text width=95mm,align=left,text=Blue,font=\fontsize{10}{12.2}\selectfont\bfseries] at (204,-62) {SUPUESTOS DEL CÁLCULO};
\node[anchor=north west,inner sep=0pt,text width=95mm,align=left,text=Navy,font=\fontsize{14}{17.1}\selectfont] at (204,-76.0) {Referencia fija: 10 g/dL};
\node[anchor=north west,inner sep=0pt,text width=95mm,align=left,text=Navy,font=\fontsize{14}{17.1}\selectfont] at (204,-90.5) {DE esperada: 1 g/dL};
\node[anchor=north west,inner sep=0pt,text width=95mm,align=left,text=Navy,font=\fontsize{14}{17.1}\selectfont\bfseries] at (204,-105.0) {Alfa: 0,05 · Potencia: 80 \%};
\node[anchor=north west,inner sep=0pt,text width=95mm,align=left,text=Navy,font=\fontsize{14}{17.1}\selectfont] at (204,-119.5) {Alternativa: mayor que 10};
\fill[Ice,rounded corners=3mm] (14,-151) rectangle (306,-165);
\node[anchor=north west,inner sep=0pt,text width=282mm,align=left,text=Navy,font=\fontsize{12}{14}\selectfont] at (19,-155) {Seguimos comparando una muestra con una referencia. Cambia la dirección del contraste.};
\node[anchor=north west,inner sep=0pt,text width=292mm,align=left,text=Gray,font=\fontsize{8.5}{10.4}\selectfont] at (14,-45) {Teoría 01, p. 7 · Teoría 02, pp. 9 y 11 · Teoría 03, pp. 18 y 19};
\end{scope}\end{tikzpicture}\mbox{}\newpage
\noindent\begin{tikzpicture}[remember picture,overlay,x=1mm,y=1mm]\begin{scope}[shift={(current page.north west)}]
\fill[Pale,rounded corners=0mm] (0,-0) rectangle (320,-180);
\node[anchor=north west,inner sep=0pt,text width=265mm,align=left,text=Blue,font=\fontsize{10}{12.2}\selectfont\bfseries] at (14,-10) {CALCULAR E INTERPRETAR  /  PUNO · VARIANTE B · INCREMENTO};
\node[anchor=north west,inner sep=0pt,text width=24mm,align=left,text=Navy,font=\fontsize{10}{12.2}\selectfont\bfseries] at (286,-10) {05 / 20};
\draw[Line,line width=.4pt] (14,-169) -- (306,-169);
\node[anchor=north west,inner sep=0pt,text width=250mm,align=left,text=Gray,font=\fontsize{9}{11.0}\selectfont] at (14,-172) {Jason Chambe  ·  Bioestadística  ·  TAREA 2};
\node[anchor=north west,inner sep=0pt,text width=34mm,align=left,text=Gray,font=\fontsize{9}{11.0}\selectfont\bfseries] at (274,-172) {DESLIZA >};
\node[anchor=north west,inner sep=0pt,text width=292mm,align=left,text=Navy,font=\fontsize{28}{34.2}\selectfont\bfseries] at (14,-25) {Del código al tamaño que reportamos};
\node[anchor=north west,inner sep=0pt,text width=292mm,align=left,text=Blue,font=\fontsize{15}{18.3}\selectfont] at (14,-44) {$d=4\qquad \mathrm{alternative}=\texttt{greater}$};
\fill[Navy,rounded corners=3mm] (14,-58) rectangle (194,-147);
\fill[white,rounded corners=3mm] (200,-58) rectangle (306,-147);
\node[anchor=north west,inner sep=0pt,text width=163mm,align=left,text=Ice,font=\fontsize{10}{12.2}\selectfont\bfseries] at (21,-64) {R · PAQUETE pwr};
\node[anchor=north west,inner sep=0pt,text width=166mm,align=left,text=white,font=\fontsize{13.5}{16.9}\selectfont] at (21,-76) {\ttfamily library(pwr)\\media\_referencia\ <-\ 10\\media\_esperada\ <-\ 14\\de\_esperada\ <-\ 1\\d\ <-\ (media\_esperada\ -\ media\_referencia)\ /\ de\_esperada\\calculo\ <-\ pwr.t.test(d\ =\ d,\ sig.level\ =\ 0.05,\\\ \ power\ =\ 0.80,\ type\ =\ "one.sample",\\\ \ alternative\ =\ "greater")\\calculo\$n\\ceiling(calculo\$n)};
\node[anchor=north west,inner sep=0pt,text width=90mm,align=left,text=Blue,font=\fontsize{10}{12.2}\selectfont\bfseries] at (208,-64) {RESULTADO};
\node[anchor=north west,inner sep=0pt,text width=91mm,align=left,text=Blue,font=\fontsize{51}{62.2}\selectfont\bfseries] at (207,-72) {3};
\node[anchor=north west,inner sep=0pt,text width=90mm,align=left,text=Navy,font=\fontsize{14}{17.1}\selectfont\bfseries] at (208,-95) {niños en una muestra};
\node[anchor=north west,inner sep=0pt,text width=90mm,align=left,text=Gray,font=\fontsize{11.5}{14.0}\selectfont] at (208,-108) {n calculado: 2,248666};
\node[anchor=north west,inner sep=0pt,text width=90mm,align=left,text=Navy,font=\fontsize{12}{14.6}\selectfont\bfseries] at (208,-120) {Redondeo hacia arriba};
\node[anchor=north west,inner sep=0pt,text width=90mm,align=left,text=Gray,font=\fontsize{11}{13.4}\selectfont] at (208,-134) {Potencia con n: 99,06 \%};
\fill[Ice,rounded corners=3mm] (14,-151) rectangle (306,-165);
\node[anchor=north west,inner sep=0pt,text width=282mm,align=left,text=Navy,font=\fontsize{12}{14}\selectfont] at (19,-154.5) {El n sin redondear disminuye, pero ambas variantes terminan en 3. Una cola no siempre cambia el entero.};
\end{scope}\end{tikzpicture}\mbox{}\newpage
\noindent\begin{tikzpicture}[remember picture,overlay,x=1mm,y=1mm]\begin{scope}[shift={(current page.north west)}]
\fill[Pale,rounded corners=0mm] (0,-0) rectangle (320,-180);
\node[anchor=north west,inner sep=0pt,text width=265mm,align=left,text=Blue,font=\fontsize{10}{12.2}\selectfont\bfseries] at (14,-10) {PLANTEAR  /  PUNO · VARIANTE C · MENOR INCREMENTO};
\node[anchor=north west,inner sep=0pt,text width=24mm,align=left,text=Navy,font=\fontsize{10}{12.2}\selectfont\bfseries] at (286,-10) {06 / 20};
\draw[Line,line width=.4pt] (14,-169) -- (306,-169);
\node[anchor=north west,inner sep=0pt,text width=250mm,align=left,text=Gray,font=\fontsize{9}{11.0}\selectfont] at (14,-172) {Jason Chambe  ·  Bioestadística  ·  TAREA 2};
\node[anchor=north west,inner sep=0pt,text width=34mm,align=left,text=Gray,font=\fontsize{9}{11.0}\selectfont\bfseries] at (274,-172) {DESLIZA >};
\node[anchor=north west,inner sep=0pt,text width=290mm,align=left,text=Navy,font=\fontsize{27}{32.9}\selectfont\bfseries] at (14,-25) {Una diferencia menor necesita más muestra};
\fill[white,rounded corners=3mm] (14,-55) rectangle (191,-147);
\fill[white,rounded corners=3mm] (197,-55) rectangle (306,-147);
\node[anchor=north west,inner sep=0pt,text width=160mm,align=left,text=Blue,font=\fontsize{10}{12.2}\selectfont\bfseries] at (21,-62) {LA COMPARACIÓN};
\node[anchor=north west,inner sep=0pt,text width=160mm,align=left,text=Navy,font=\fontsize{18}{23}\selectfont] at (21,-73) {Ahora esperamos 11 g/dL en lugar de 14. Mantenemos la referencia, la DE, el alfa y la potencia del caso bilateral.};
\node[anchor=north west,inner sep=0pt,text width=160mm,align=left,text=Blue,font=\fontsize{25}{30.5}\selectfont\bfseries] at (21,-108) {$10\ \mathrm{g/dL}\ \longrightarrow\ 11\ \mathrm{g/dL}$};
\node[anchor=north west,inner sep=0pt,text width=160mm,align=left,text=Navy,font=\fontsize{12.4}{15.1}\selectfont] at (21,-129) {$H_0:\mu=10\qquad H_A:\mu\ne10$};
\node[anchor=north west,inner sep=0pt,text width=95mm,align=left,text=Blue,font=\fontsize{10}{12.2}\selectfont\bfseries] at (204,-62) {SUPUESTOS DEL CÁLCULO};
\node[anchor=north west,inner sep=0pt,text width=95mm,align=left,text=Navy,font=\fontsize{14}{17.1}\selectfont] at (204,-76.0) {Referencia fija: 10 g/dL};
\node[anchor=north west,inner sep=0pt,text width=95mm,align=left,text=Navy,font=\fontsize{14}{17.1}\selectfont] at (204,-90.5) {DE esperada: 1 g/dL};
\node[anchor=north west,inner sep=0pt,text width=95mm,align=left,text=Navy,font=\fontsize{14}{17.1}\selectfont\bfseries] at (204,-105.0) {Alfa: 0,05 · Potencia: 80 \%};
\node[anchor=north west,inner sep=0pt,text width=95mm,align=left,text=Navy,font=\fontsize{14}{17.1}\selectfont] at (204,-119.5) {Alternativa: bilateral};
\fill[Ice,rounded corners=3mm] (14,-151) rectangle (306,-165);
\node[anchor=north west,inner sep=0pt,text width=282mm,align=left,text=Navy,font=\fontsize{12}{14}\selectfont] at (19,-155) {La comparación sigue siendo de una muestra. Cambia la diferencia que queremos detectar.};
\node[anchor=north west,inner sep=0pt,text width=292mm,align=left,text=Gray,font=\fontsize{8.5}{10.4}\selectfont] at (14,-45) {Teoría 01, p. 7 · Teoría 02, pp. 9 y 11 · Teoría 03, pp. 15 y 19};
\end{scope}\end{tikzpicture}\mbox{}\newpage
\noindent\begin{tikzpicture}[remember picture,overlay,x=1mm,y=1mm]\begin{scope}[shift={(current page.north west)}]
\fill[Pale,rounded corners=0mm] (0,-0) rectangle (320,-180);
\node[anchor=north west,inner sep=0pt,text width=265mm,align=left,text=Blue,font=\fontsize{10}{12.2}\selectfont\bfseries] at (14,-10) {CALCULAR E INTERPRETAR  /  PUNO · VARIANTE C · MENOR INCREMENTO};
\node[anchor=north west,inner sep=0pt,text width=24mm,align=left,text=Navy,font=\fontsize{10}{12.2}\selectfont\bfseries] at (286,-10) {07 / 20};
\draw[Line,line width=.4pt] (14,-169) -- (306,-169);
\node[anchor=north west,inner sep=0pt,text width=250mm,align=left,text=Gray,font=\fontsize{9}{11.0}\selectfont] at (14,-172) {Jason Chambe  ·  Bioestadística  ·  TAREA 2};
\node[anchor=north west,inner sep=0pt,text width=34mm,align=left,text=Gray,font=\fontsize{9}{11.0}\selectfont\bfseries] at (274,-172) {DESLIZA >};
\node[anchor=north west,inner sep=0pt,text width=292mm,align=left,text=Navy,font=\fontsize{28}{34.2}\selectfont\bfseries] at (14,-25) {Del código al tamaño que reportamos};
\node[anchor=north west,inner sep=0pt,text width=292mm,align=left,text=Blue,font=\fontsize{15}{18.3}\selectfont] at (14,-44) {$d=(11-10)/1=1$};
\fill[Navy,rounded corners=3mm] (14,-58) rectangle (194,-147);
\fill[white,rounded corners=3mm] (200,-58) rectangle (306,-147);
\node[anchor=north west,inner sep=0pt,text width=163mm,align=left,text=Ice,font=\fontsize{10}{12.2}\selectfont\bfseries] at (21,-64) {R · PAQUETE pwr};
\node[anchor=north west,inner sep=0pt,text width=166mm,align=left,text=white,font=\fontsize{13.5}{16.9}\selectfont] at (21,-76) {\ttfamily library(pwr)\\media\_referencia\ <-\ 10\\media\_esperada\ <-\ 11\\de\_esperada\ <-\ 1\\d\ <-\ (media\_esperada\ -\ media\_referencia)\ /\ de\_esperada\\calculo\ <-\ pwr.t.test(d\ =\ d,\ sig.level\ =\ 0.05,\\\ \ power\ =\ 0.80,\ type\ =\ "one.sample",\\\ \ alternative\ =\ "two.sided")\\calculo\$n\\ceiling(calculo\$n)};
\node[anchor=north west,inner sep=0pt,text width=90mm,align=left,text=Blue,font=\fontsize{10}{12.2}\selectfont\bfseries] at (208,-64) {RESULTADO};
\node[anchor=north west,inner sep=0pt,text width=91mm,align=left,text=Blue,font=\fontsize{51}{62.2}\selectfont\bfseries] at (207,-72) {10};
\node[anchor=north west,inner sep=0pt,text width=90mm,align=left,text=Navy,font=\fontsize{14}{17.1}\selectfont\bfseries] at (208,-95) {niños en una muestra};
\node[anchor=north west,inner sep=0pt,text width=90mm,align=left,text=Gray,font=\fontsize{11.5}{14.0}\selectfont] at (208,-108) {n calculado: 9,937850};
\node[anchor=north west,inner sep=0pt,text width=90mm,align=left,text=Navy,font=\fontsize{12}{14.6}\selectfont\bfseries] at (208,-120) {Redondeo hacia arriba};
\node[anchor=north west,inner sep=0pt,text width=90mm,align=left,text=Gray,font=\fontsize{11}{13.4}\selectfont] at (208,-134) {Potencia con n: 80,31 \%};
\fill[Ice,rounded corners=3mm] (14,-151) rectangle (306,-165);
\node[anchor=north west,inner sep=0pt,text width=282mm,align=left,text=Navy,font=\fontsize{12}{14}\selectfont] at (19,-154.5) {El efecto baja de d = 4 a d = 1. Con los demás supuestos iguales, la muestra aumenta de 3 a 10.};
\end{scope}\end{tikzpicture}\mbox{}\newpage
\noindent\begin{tikzpicture}[remember picture,overlay,x=1mm,y=1mm]\begin{scope}[shift={(current page.north west)}]
\fill[Pale,rounded corners=0mm] (0,-0) rectangle (320,-180);
\node[anchor=north west,inner sep=0pt,text width=265mm,align=left,text=Blue,font=\fontsize{10}{12.2}\selectfont\bfseries] at (14,-10) {PLANTEAR  /  JULIACA · VARIANTE A · BILATERAL};
\node[anchor=north west,inner sep=0pt,text width=24mm,align=left,text=Navy,font=\fontsize{10}{12.2}\selectfont\bfseries] at (286,-10) {08 / 20};
\draw[Line,line width=.4pt] (14,-169) -- (306,-169);
\node[anchor=north west,inner sep=0pt,text width=250mm,align=left,text=Gray,font=\fontsize{9}{11.0}\selectfont] at (14,-172) {Jason Chambe  ·  Bioestadística  ·  TAREA 2};
\node[anchor=north west,inner sep=0pt,text width=34mm,align=left,text=Gray,font=\fontsize{9}{11.0}\selectfont\bfseries] at (274,-172) {DESLIZA >};
\node[anchor=north west,inner sep=0pt,text width=290mm,align=left,text=Navy,font=\fontsize{27}{32.9}\selectfont\bfseries] at (14,-25) {Anemia: una proporción frente a referencia};
\fill[white,rounded corners=3mm] (14,-55) rectangle (191,-147);
\fill[white,rounded corners=3mm] (197,-55) rectangle (306,-147);
\node[anchor=north west,inner sep=0pt,text width=160mm,align=left,text=Blue,font=\fontsize{10}{12.2}\selectfont\bfseries] at (21,-62) {LA COMPARACIÓN};
\node[anchor=north west,inner sep=0pt,text width=160mm,align=left,text=Navy,font=\fontsize{18}{23}\selectfont] at (21,-73) {Esperamos una prevalencia posterior de 20 \%. Para este cálculo tratamos el 40 \% previo como referencia fija.};
\node[anchor=north west,inner sep=0pt,text width=160mm,align=left,text=Blue,font=\fontsize{25}{30.5}\selectfont\bfseries] at (21,-108) {$40\,\%\ \longrightarrow\ 20\,\%$};
\node[anchor=north west,inner sep=0pt,text width=160mm,align=left,text=Navy,font=\fontsize{12.4}{15.1}\selectfont] at (21,-129) {$H_0:p=0{,}40\qquad H_A:p\ne0{,}40$};
\node[anchor=north west,inner sep=0pt,text width=95mm,align=left,text=Blue,font=\fontsize{10}{12.2}\selectfont\bfseries] at (204,-62) {SUPUESTOS DEL CÁLCULO};
\node[anchor=north west,inner sep=0pt,text width=95mm,align=left,text=Navy,font=\fontsize{14}{17.1}\selectfont] at (204,-76.0) {Desenlace: anemia sí/no};
\node[anchor=north west,inner sep=0pt,text width=95mm,align=left,text=Navy,font=\fontsize{14}{17.1}\selectfont] at (204,-90.5) {Referencia fija: 40 \%};
\node[anchor=north west,inner sep=0pt,text width=95mm,align=left,text=Navy,font=\fontsize{14}{17.1}\selectfont\bfseries] at (204,-105.0) {Alfa: 0,05 · Potencia: 80 \%};
\node[anchor=north west,inner sep=0pt,text width=95mm,align=left,text=Navy,font=\fontsize{14}{17.1}\selectfont] at (204,-119.5) {Alternativa: bilateral};
\fill[Ice,rounded corners=3mm] (14,-151) rectangle (306,-165);
\node[anchor=north west,inner sep=0pt,text width=282mm,align=left,text=Navy,font=\fontsize{12}{14}\selectfont] at (19,-155) {Se mide prevalencia en una muestra nueva. Una comparación histórica no demuestra por sí sola causalidad.};
\node[anchor=north west,inner sep=0pt,text width=292mm,align=left,text=Gray,font=\fontsize{8.5}{10.4}\selectfont] at (14,-45) {Teoría 01, pp. 47 y 50 · Teoría 02, pp. 9 y 11 · Teoría 03, pp. 32 y 33};
\end{scope}\end{tikzpicture}\mbox{}\newpage
\noindent\begin{tikzpicture}[remember picture,overlay,x=1mm,y=1mm]\begin{scope}[shift={(current page.north west)}]
\fill[Pale,rounded corners=0mm] (0,-0) rectangle (320,-180);
\node[anchor=north west,inner sep=0pt,text width=265mm,align=left,text=Blue,font=\fontsize{10}{12.2}\selectfont\bfseries] at (14,-10) {CALCULAR E INTERPRETAR  /  JULIACA · VARIANTE A · BILATERAL};
\node[anchor=north west,inner sep=0pt,text width=24mm,align=left,text=Navy,font=\fontsize{10}{12.2}\selectfont\bfseries] at (286,-10) {09 / 20};
\draw[Line,line width=.4pt] (14,-169) -- (306,-169);
\node[anchor=north west,inner sep=0pt,text width=250mm,align=left,text=Gray,font=\fontsize{9}{11.0}\selectfont] at (14,-172) {Jason Chambe  ·  Bioestadística  ·  TAREA 2};
\node[anchor=north west,inner sep=0pt,text width=34mm,align=left,text=Gray,font=\fontsize{9}{11.0}\selectfont\bfseries] at (274,-172) {DESLIZA >};
\node[anchor=north west,inner sep=0pt,text width=292mm,align=left,text=Navy,font=\fontsize{28}{34.2}\selectfont\bfseries] at (14,-25) {Del código al tamaño que reportamos};
\node[anchor=north west,inner sep=0pt,text width=292mm,align=left,text=Blue,font=\fontsize{15}{18.3}\selectfont] at (14,-44) {$h=2\arcsin\sqrt{0{,}20}-2\arcsin\sqrt{0{,}40}$};
\fill[Navy,rounded corners=3mm] (14,-58) rectangle (194,-147);
\fill[white,rounded corners=3mm] (200,-58) rectangle (306,-147);
\node[anchor=north west,inner sep=0pt,text width=163mm,align=left,text=Ice,font=\fontsize{10}{12.2}\selectfont\bfseries] at (21,-64) {R · PAQUETE pwr};
\node[anchor=north west,inner sep=0pt,text width=166mm,align=left,text=white,font=\fontsize{13.5}{16.9}\selectfont] at (21,-76) {\ttfamily library(pwr)\\p\_referencia\ <-\ 0.40\\p\_esperada\ <-\ 0.20\\h\ <-\ ES.h(p\_esperada,\ p\_referencia)\\calculo\ <-\ pwr.p.test(h\ =\ h,\ sig.level\ =\ 0.05,\\\ \ power\ =\ 0.80,\ alternative\ =\ "two.sided")\\h\\calculo\$n\\ceiling(calculo\$n)};
\node[anchor=north west,inner sep=0pt,text width=90mm,align=left,text=Blue,font=\fontsize{10}{12.2}\selectfont\bfseries] at (208,-64) {RESULTADO};
\node[anchor=north west,inner sep=0pt,text width=91mm,align=left,text=Blue,font=\fontsize{51}{62.2}\selectfont\bfseries] at (207,-72) {41};
\node[anchor=north west,inner sep=0pt,text width=90mm,align=left,text=Navy,font=\fontsize{14}{17.1}\selectfont\bfseries] at (208,-95) {niños en una muestra};
\node[anchor=north west,inner sep=0pt,text width=90mm,align=left,text=Gray,font=\fontsize{11.5}{14.0}\selectfont] at (208,-108) {n calculado: 40,149573};
\node[anchor=north west,inner sep=0pt,text width=90mm,align=left,text=Navy,font=\fontsize{12}{14.6}\selectfont\bfseries] at (208,-120) {Redondeo hacia arriba};
\node[anchor=north west,inner sep=0pt,text width=90mm,align=left,text=Gray,font=\fontsize{11}{13.4}\selectfont] at (208,-134) {Potencia con n: 80,82 \%};
\fill[Ice,rounded corners=3mm] (14,-151) rectangle (306,-165);
\node[anchor=north west,inner sep=0pt,text width=282mm,align=left,text=Navy,font=\fontsize{12}{14}\selectfont] at (19,-154.5) {40 \% a 20 \% son 20 puntos porcentuales de reducción. El n corresponde a una sola muestra.};
\end{scope}\end{tikzpicture}\mbox{}\newpage
\noindent\begin{tikzpicture}[remember picture,overlay,x=1mm,y=1mm]\begin{scope}[shift={(current page.north west)}]
\fill[Pale,rounded corners=0mm] (0,-0) rectangle (320,-180);
\node[anchor=north west,inner sep=0pt,text width=265mm,align=left,text=Blue,font=\fontsize{10}{12.2}\selectfont\bfseries] at (14,-10) {PLANTEAR  /  JULIACA · VARIANTE B · UNILATERAL};
\node[anchor=north west,inner sep=0pt,text width=24mm,align=left,text=Navy,font=\fontsize{10}{12.2}\selectfont\bfseries] at (286,-10) {10 / 20};
\draw[Line,line width=.4pt] (14,-169) -- (306,-169);
\node[anchor=north west,inner sep=0pt,text width=250mm,align=left,text=Gray,font=\fontsize{9}{11.0}\selectfont] at (14,-172) {Jason Chambe  ·  Bioestadística  ·  TAREA 2};
\node[anchor=north west,inner sep=0pt,text width=34mm,align=left,text=Gray,font=\fontsize{9}{11.0}\selectfont\bfseries] at (274,-172) {DESLIZA >};
\node[anchor=north west,inner sep=0pt,text width=290mm,align=left,text=Navy,font=\fontsize{27}{32.9}\selectfont\bfseries] at (14,-25) {Una hipótesis de disminución};
\fill[white,rounded corners=3mm] (14,-55) rectangle (191,-147);
\fill[white,rounded corners=3mm] (197,-55) rectangle (306,-147);
\node[anchor=north west,inner sep=0pt,text width=160mm,align=left,text=Blue,font=\fontsize{10}{12.2}\selectfont\bfseries] at (21,-62) {LA COMPARACIÓN};
\node[anchor=north west,inner sep=0pt,text width=160mm,align=left,text=Navy,font=\fontsize{18}{23}\selectfont] at (21,-73) {Se mantienen 40 \% y 20 \%, pero se plantea antes del estudio que el contraste evaluará una disminución.};
\node[anchor=north west,inner sep=0pt,text width=160mm,align=left,text=Blue,font=\fontsize{25}{30.5}\selectfont\bfseries] at (21,-108) {$40\,\%\ \longrightarrow\ 20\,\%$};
\node[anchor=north west,inner sep=0pt,text width=160mm,align=left,text=Navy,font=\fontsize{12.4}{15.1}\selectfont] at (21,-129) {$H_0:p\ge0{,}40\qquad H_A:p<0{,}40$};
\node[anchor=north west,inner sep=0pt,text width=95mm,align=left,text=Blue,font=\fontsize{10}{12.2}\selectfont\bfseries] at (204,-62) {SUPUESTOS DEL CÁLCULO};
\node[anchor=north west,inner sep=0pt,text width=95mm,align=left,text=Navy,font=\fontsize{14}{17.1}\selectfont] at (204,-76.0) {Desenlace: anemia sí/no};
\node[anchor=north west,inner sep=0pt,text width=95mm,align=left,text=Navy,font=\fontsize{14}{17.1}\selectfont] at (204,-90.5) {Referencia fija: 40 \%};
\node[anchor=north west,inner sep=0pt,text width=95mm,align=left,text=Navy,font=\fontsize{14}{17.1}\selectfont\bfseries] at (204,-105.0) {Alfa: 0,05 · Potencia: 80 \%};
\node[anchor=north west,inner sep=0pt,text width=95mm,align=left,text=Navy,font=\fontsize{14}{17.1}\selectfont] at (204,-119.5) {Alternativa: menor que 40 \%};
\fill[Ice,rounded corners=3mm] (14,-151) rectangle (306,-165);
\node[anchor=north west,inner sep=0pt,text width=282mm,align=left,text=Navy,font=\fontsize{12}{14}\selectfont] at (19,-155) {La comparación sigue siendo de una proporción. La hipótesis direccional debe estar justificada de antemano.};
\node[anchor=north west,inner sep=0pt,text width=292mm,align=left,text=Gray,font=\fontsize{8.5}{10.4}\selectfont] at (14,-45) {Teoría 01, p. 47 · Teoría 02, pp. 9 y 11 · Teoría 03, pp. 18 y 19};
\end{scope}\end{tikzpicture}\mbox{}\newpage
\noindent\begin{tikzpicture}[remember picture,overlay,x=1mm,y=1mm]\begin{scope}[shift={(current page.north west)}]
\fill[Pale,rounded corners=0mm] (0,-0) rectangle (320,-180);
\node[anchor=north west,inner sep=0pt,text width=265mm,align=left,text=Blue,font=\fontsize{10}{12.2}\selectfont\bfseries] at (14,-10) {CALCULAR E INTERPRETAR  /  JULIACA · VARIANTE B · UNILATERAL};
\node[anchor=north west,inner sep=0pt,text width=24mm,align=left,text=Navy,font=\fontsize{10}{12.2}\selectfont\bfseries] at (286,-10) {11 / 20};
\draw[Line,line width=.4pt] (14,-169) -- (306,-169);
\node[anchor=north west,inner sep=0pt,text width=250mm,align=left,text=Gray,font=\fontsize{9}{11.0}\selectfont] at (14,-172) {Jason Chambe  ·  Bioestadística  ·  TAREA 2};
\node[anchor=north west,inner sep=0pt,text width=34mm,align=left,text=Gray,font=\fontsize{9}{11.0}\selectfont\bfseries] at (274,-172) {DESLIZA >};
\node[anchor=north west,inner sep=0pt,text width=292mm,align=left,text=Navy,font=\fontsize{28}{34.2}\selectfont\bfseries] at (14,-25) {Del código al tamaño que reportamos};
\node[anchor=north west,inner sep=0pt,text width=292mm,align=left,text=Blue,font=\fontsize{15}{18.3}\selectfont] at (14,-44) {$h=-0{,}442143\qquad\mathrm{alternative}=\texttt{less}$};
\fill[Navy,rounded corners=3mm] (14,-58) rectangle (194,-147);
\fill[white,rounded corners=3mm] (200,-58) rectangle (306,-147);
\node[anchor=north west,inner sep=0pt,text width=163mm,align=left,text=Ice,font=\fontsize{10}{12.2}\selectfont\bfseries] at (21,-64) {R · PAQUETE pwr};
\node[anchor=north west,inner sep=0pt,text width=166mm,align=left,text=white,font=\fontsize{13.5}{16.9}\selectfont] at (21,-76) {\ttfamily library(pwr)\\p\_referencia\ <-\ 0.40\\p\_esperada\ <-\ 0.20\\h\ <-\ ES.h(p\_esperada,\ p\_referencia)\\calculo\ <-\ pwr.p.test(h\ =\ h,\ sig.level\ =\ 0.05,\\\ \ power\ =\ 0.80,\ alternative\ =\ "less")\\h\\calculo\$n\\ceiling(calculo\$n)};
\node[anchor=north west,inner sep=0pt,text width=90mm,align=left,text=Blue,font=\fontsize{10}{12.2}\selectfont\bfseries] at (208,-64) {RESULTADO};
\node[anchor=north west,inner sep=0pt,text width=91mm,align=left,text=Blue,font=\fontsize{51}{62.2}\selectfont\bfseries] at (207,-72) {32};
\node[anchor=north west,inner sep=0pt,text width=90mm,align=left,text=Navy,font=\fontsize{14}{17.1}\selectfont\bfseries] at (208,-95) {niños en una muestra};
\node[anchor=north west,inner sep=0pt,text width=90mm,align=left,text=Gray,font=\fontsize{11.5}{14.0}\selectfont] at (208,-108) {n calculado: 31,625855};
\node[anchor=north west,inner sep=0pt,text width=90mm,align=left,text=Navy,font=\fontsize{12}{14.6}\selectfont\bfseries] at (208,-120) {Redondeo hacia arriba};
\node[anchor=north west,inner sep=0pt,text width=90mm,align=left,text=Gray,font=\fontsize{11}{13.4}\selectfont] at (208,-134) {Potencia con n: 80,41 \%};
\fill[Ice,rounded corners=3mm] (14,-151) rectangle (306,-165);
\node[anchor=north west,inner sep=0pt,text width=282mm,align=left,text=Navy,font=\fontsize{12}{14}\selectfont] at (19,-154.5) {Conserva h negativo y usa less. Cambiar a una cola después de ver los datos no es una justificación válida.};
\end{scope}\end{tikzpicture}\mbox{}\newpage
\noindent\begin{tikzpicture}[remember picture,overlay,x=1mm,y=1mm]\begin{scope}[shift={(current page.north west)}]
\fill[Pale,rounded corners=0mm] (0,-0) rectangle (320,-180);
\node[anchor=north west,inner sep=0pt,text width=265mm,align=left,text=Blue,font=\fontsize{10}{12.2}\selectfont\bfseries] at (14,-10) {PLANTEAR  /  JULIACA · VARIANTE C · MENOR DIFERENCIA};
\node[anchor=north west,inner sep=0pt,text width=24mm,align=left,text=Navy,font=\fontsize{10}{12.2}\selectfont\bfseries] at (286,-10) {12 / 20};
\draw[Line,line width=.4pt] (14,-169) -- (306,-169);
\node[anchor=north west,inner sep=0pt,text width=250mm,align=left,text=Gray,font=\fontsize{9}{11.0}\selectfont] at (14,-172) {Jason Chambe  ·  Bioestadística  ·  TAREA 2};
\node[anchor=north west,inner sep=0pt,text width=34mm,align=left,text=Gray,font=\fontsize{9}{11.0}\selectfont\bfseries] at (274,-172) {DESLIZA >};
\node[anchor=north west,inner sep=0pt,text width=290mm,align=left,text=Navy,font=\fontsize{27}{32.9}\selectfont\bfseries] at (14,-25) {Detectar una reducción más discreta};
\fill[white,rounded corners=3mm] (14,-55) rectangle (191,-147);
\fill[white,rounded corners=3mm] (197,-55) rectangle (306,-147);
\node[anchor=north west,inner sep=0pt,text width=160mm,align=left,text=Blue,font=\fontsize{10}{12.2}\selectfont\bfseries] at (21,-62) {LA COMPARACIÓN};
\node[anchor=north west,inner sep=0pt,text width=160mm,align=left,text=Navy,font=\fontsize{18}{23}\selectfont] at (21,-73) {Ahora se espera pasar de 40 \% a 30 \%. Conservamos el contraste bilateral, alfa de 0,05 y potencia de 80 \%.};
\node[anchor=north west,inner sep=0pt,text width=160mm,align=left,text=Blue,font=\fontsize{25}{30.5}\selectfont\bfseries] at (21,-108) {$40\,\%\ \longrightarrow\ 30\,\%$};
\node[anchor=north west,inner sep=0pt,text width=160mm,align=left,text=Navy,font=\fontsize{12.4}{15.1}\selectfont] at (21,-129) {$H_0:p=0{,}40\qquad H_A:p\ne0{,}40$};
\node[anchor=north west,inner sep=0pt,text width=95mm,align=left,text=Blue,font=\fontsize{10}{12.2}\selectfont\bfseries] at (204,-62) {SUPUESTOS DEL CÁLCULO};
\node[anchor=north west,inner sep=0pt,text width=95mm,align=left,text=Navy,font=\fontsize{14}{17.1}\selectfont] at (204,-76.0) {Desenlace: anemia sí/no};
\node[anchor=north west,inner sep=0pt,text width=95mm,align=left,text=Navy,font=\fontsize{14}{17.1}\selectfont] at (204,-90.5) {Referencia fija: 40 \%};
\node[anchor=north west,inner sep=0pt,text width=95mm,align=left,text=Navy,font=\fontsize{14}{17.1}\selectfont\bfseries] at (204,-105.0) {Alfa: 0,05 · Potencia: 80 \%};
\node[anchor=north west,inner sep=0pt,text width=95mm,align=left,text=Navy,font=\fontsize{14}{17.1}\selectfont] at (204,-119.5) {Alternativa: bilateral};
\fill[Ice,rounded corners=3mm] (14,-151) rectangle (306,-165);
\node[anchor=north west,inner sep=0pt,text width=282mm,align=left,text=Navy,font=\fontsize{12}{14}\selectfont] at (19,-155) {Mantenemos una muestra contra referencia. Solo cambia la prevalencia que esperamos detectar.};
\node[anchor=north west,inner sep=0pt,text width=292mm,align=left,text=Gray,font=\fontsize{8.5}{10.4}\selectfont] at (14,-45) {Teoría 01, pp. 47 y 50 · Teoría 02, p. 11 · Teoría 03, pp. 15 y 19};
\end{scope}\end{tikzpicture}\mbox{}\newpage
\noindent\begin{tikzpicture}[remember picture,overlay,x=1mm,y=1mm]\begin{scope}[shift={(current page.north west)}]
\fill[Pale,rounded corners=0mm] (0,-0) rectangle (320,-180);
\node[anchor=north west,inner sep=0pt,text width=265mm,align=left,text=Blue,font=\fontsize{10}{12.2}\selectfont\bfseries] at (14,-10) {CALCULAR E INTERPRETAR  /  JULIACA · VARIANTE C · MENOR DIFERENCIA};
\node[anchor=north west,inner sep=0pt,text width=24mm,align=left,text=Navy,font=\fontsize{10}{12.2}\selectfont\bfseries] at (286,-10) {13 / 20};
\draw[Line,line width=.4pt] (14,-169) -- (306,-169);
\node[anchor=north west,inner sep=0pt,text width=250mm,align=left,text=Gray,font=\fontsize{9}{11.0}\selectfont] at (14,-172) {Jason Chambe  ·  Bioestadística  ·  TAREA 2};
\node[anchor=north west,inner sep=0pt,text width=34mm,align=left,text=Gray,font=\fontsize{9}{11.0}\selectfont\bfseries] at (274,-172) {DESLIZA >};
\node[anchor=north west,inner sep=0pt,text width=292mm,align=left,text=Navy,font=\fontsize{28}{34.2}\selectfont\bfseries] at (14,-25) {Del código al tamaño que reportamos};
\node[anchor=north west,inner sep=0pt,text width=292mm,align=left,text=Blue,font=\fontsize{15}{18.3}\selectfont] at (14,-44) {$h=2\arcsin\sqrt{0{,}30}-2\arcsin\sqrt{0{,}40}$};
\fill[Navy,rounded corners=3mm] (14,-58) rectangle (194,-147);
\fill[white,rounded corners=3mm] (200,-58) rectangle (306,-147);
\node[anchor=north west,inner sep=0pt,text width=163mm,align=left,text=Ice,font=\fontsize{10}{12.2}\selectfont\bfseries] at (21,-64) {R · PAQUETE pwr};
\node[anchor=north west,inner sep=0pt,text width=166mm,align=left,text=white,font=\fontsize{13.5}{16.9}\selectfont] at (21,-76) {\ttfamily library(pwr)\\p\_referencia\ <-\ 0.40\\p\_esperada\ <-\ 0.30\\h\ <-\ ES.h(p\_esperada,\ p\_referencia)\\calculo\ <-\ pwr.p.test(h\ =\ h,\ sig.level\ =\ 0.05,\\\ \ power\ =\ 0.80,\ alternative\ =\ "two.sided")\\h\\calculo\$n\\ceiling(calculo\$n)};
\node[anchor=north west,inner sep=0pt,text width=90mm,align=left,text=Blue,font=\fontsize{10}{12.2}\selectfont\bfseries] at (208,-64) {RESULTADO};
\node[anchor=north west,inner sep=0pt,text width=91mm,align=left,text=Blue,font=\fontsize{51}{62.2}\selectfont\bfseries] at (207,-72) {178};
\node[anchor=north west,inner sep=0pt,text width=90mm,align=left,text=Navy,font=\fontsize{14}{17.1}\selectfont\bfseries] at (208,-95) {niños en una muestra};
\node[anchor=north west,inner sep=0pt,text width=90mm,align=left,text=Gray,font=\fontsize{11.5}{14.0}\selectfont] at (208,-108) {n calculado: 177,709631};
\node[anchor=north west,inner sep=0pt,text width=90mm,align=left,text=Navy,font=\fontsize{12}{14.6}\selectfont\bfseries] at (208,-120) {Redondeo hacia arriba};
\node[anchor=north west,inner sep=0pt,text width=90mm,align=left,text=Gray,font=\fontsize{11}{13.4}\selectfont] at (208,-134) {Potencia con n: 80,06 \%};
\fill[Ice,rounded corners=3mm] (14,-151) rectangle (306,-165);
\node[anchor=north west,inner sep=0pt,text width=282mm,align=left,text=Navy,font=\fontsize{12}{14}\selectfont] at (19,-154.5) {Al reducir la diferencia de 20 a 10 puntos porcentuales, la muestra aumenta de 41 a 178 niños.};
\end{scope}\end{tikzpicture}\mbox{}\newpage
\noindent\begin{tikzpicture}[remember picture,overlay,x=1mm,y=1mm]\begin{scope}[shift={(current page.north west)}]
\fill[Pale,rounded corners=0mm] (0,-0) rectangle (320,-180);
\node[anchor=north west,inner sep=0pt,text width=265mm,align=left,text=Blue,font=\fontsize{10}{12.2}\selectfont\bfseries] at (14,-10) {PLANTEAR  /  SJL Y VMT · VARIANTE A · POTENCIA DE 80 \%};
\node[anchor=north west,inner sep=0pt,text width=24mm,align=left,text=Navy,font=\fontsize{10}{12.2}\selectfont\bfseries] at (286,-10) {14 / 20};
\draw[Line,line width=.4pt] (14,-169) -- (306,-169);
\node[anchor=north west,inner sep=0pt,text width=250mm,align=left,text=Gray,font=\fontsize{9}{11.0}\selectfont] at (14,-172) {Jason Chambe  ·  Bioestadística  ·  TAREA 2};
\node[anchor=north west,inner sep=0pt,text width=34mm,align=left,text=Gray,font=\fontsize{9}{11.0}\selectfont\bfseries] at (274,-172) {DESLIZA >};
\node[anchor=north west,inner sep=0pt,text width=290mm,align=left,text=Navy,font=\fontsize{27}{32.9}\selectfont\bfseries] at (14,-25) {Dos distritos, dos muestras independientes};
\fill[white,rounded corners=3mm] (14,-55) rectangle (191,-147);
\fill[white,rounded corners=3mm] (197,-55) rectangle (306,-147);
\node[anchor=north west,inner sep=0pt,text width=160mm,align=left,text=Blue,font=\fontsize{10}{12.2}\selectfont\bfseries] at (21,-62) {LA COMPARACIÓN};
\node[anchor=north west,inner sep=0pt,text width=160mm,align=left,text=Navy,font=\fontsize{18}{23}\selectfont] at (21,-73) {Comparemos 43 \% de prevalencia en San Juan de Lurigancho y 26 \% en Villa María del Triunfo.};
\node[anchor=north west,inner sep=0pt,text width=160mm,align=left,text=Blue,font=\fontsize{25}{30.5}\selectfont\bfseries] at (21,-108) {$43\,\%\quad\mathrm{vs.}\quad26\,\%$};
\node[anchor=north west,inner sep=0pt,text width=160mm,align=left,text=Navy,font=\fontsize{12.4}{15.1}\selectfont] at (21,-129) {$H_0:p_{SJL}=p_{VMT}\quad H_A:p_{SJL}\ne p_{VMT}$};
\node[anchor=north west,inner sep=0pt,text width=95mm,align=left,text=Blue,font=\fontsize{10}{12.2}\selectfont\bfseries] at (204,-62) {SUPUESTOS DEL CÁLCULO};
\node[anchor=north west,inner sep=0pt,text width=95mm,align=left,text=Navy,font=\fontsize{14}{17.1}\selectfont] at (204,-76.0) {Niños distintos por distrito};
\node[anchor=north west,inner sep=0pt,text width=95mm,align=left,text=Navy,font=\fontsize{14}{17.1}\selectfont] at (204,-90.5) {Asignación de muestra: 1:1};
\node[anchor=north west,inner sep=0pt,text width=95mm,align=left,text=Navy,font=\fontsize{14}{17.1}\selectfont\bfseries] at (204,-105.0) {Alfa: 0,05 · Potencia: 80 \%};
\node[anchor=north west,inner sep=0pt,text width=95mm,align=left,text=Navy,font=\fontsize{14}{17.1}\selectfont] at (204,-119.5) {Alternativa: bilateral};
\fill[Ice,rounded corners=3mm] (14,-151) rectangle (306,-165);
\node[anchor=north west,inner sep=0pt,text width=282mm,align=left,text=Navy,font=\fontsize{12}{14}\selectfont] at (19,-155) {Es una comparación transversal observacional. Aquí ambas prevalencias se estimarán mediante muestras.};
\node[anchor=north west,inner sep=0pt,text width=292mm,align=left,text=Gray,font=\fontsize{8.5}{10.4}\selectfont] at (14,-45) {Teoría 01, pp. 29 y 47 · Teoría 02, p. 7 · Teoría 03, pp. 36 y 37};
\end{scope}\end{tikzpicture}\mbox{}\newpage
\noindent\begin{tikzpicture}[remember picture,overlay,x=1mm,y=1mm]\begin{scope}[shift={(current page.north west)}]
\fill[Pale,rounded corners=0mm] (0,-0) rectangle (320,-180);
\node[anchor=north west,inner sep=0pt,text width=265mm,align=left,text=Blue,font=\fontsize{10}{12.2}\selectfont\bfseries] at (14,-10) {CALCULAR E INTERPRETAR  /  SJL Y VMT · VARIANTE A · POTENCIA DE 80 \%};
\node[anchor=north west,inner sep=0pt,text width=24mm,align=left,text=Navy,font=\fontsize{10}{12.2}\selectfont\bfseries] at (286,-10) {15 / 20};
\draw[Line,line width=.4pt] (14,-169) -- (306,-169);
\node[anchor=north west,inner sep=0pt,text width=250mm,align=left,text=Gray,font=\fontsize{9}{11.0}\selectfont] at (14,-172) {Jason Chambe  ·  Bioestadística  ·  TAREA 2};
\node[anchor=north west,inner sep=0pt,text width=34mm,align=left,text=Gray,font=\fontsize{9}{11.0}\selectfont\bfseries] at (274,-172) {DESLIZA >};
\node[anchor=north west,inner sep=0pt,text width=292mm,align=left,text=Navy,font=\fontsize{28}{34.2}\selectfont\bfseries] at (14,-25) {Del código al tamaño que reportamos};
\node[anchor=north west,inner sep=0pt,text width=292mm,align=left,text=Blue,font=\fontsize{15}{18.3}\selectfont] at (14,-44) {$h=2\arcsin\sqrt{0{,}43}-2\arcsin\sqrt{0{,}26}$};
\fill[Navy,rounded corners=3mm] (14,-58) rectangle (194,-147);
\fill[white,rounded corners=3mm] (200,-58) rectangle (306,-147);
\node[anchor=north west,inner sep=0pt,text width=163mm,align=left,text=Ice,font=\fontsize{10}{12.2}\selectfont\bfseries] at (21,-64) {R · PAQUETE pwr};
\node[anchor=north west,inner sep=0pt,text width=166mm,align=left,text=white,font=\fontsize{13.5}{16.9}\selectfont] at (21,-76) {\ttfamily library(pwr)\\p\_sjl\ <-\ 0.43\\p\_vmt\ <-\ 0.26\\h\ <-\ ES.h(p\_sjl,\ p\_vmt)\\calculo\ <-\ pwr.2p.test(h\ =\ h,\ sig.level\ =\ 0.05,\\\ \ power\ =\ 0.80,\ alternative\ =\ "two.sided")\\n\_por\_distrito\ <-\ ceiling(calculo\$n)\\calculo\$n\\c(por\_distrito\ =\ n\_por\_distrito,\\\ \ total\ =\ 2\ *\ n\_por\_distrito)};
\node[anchor=north west,inner sep=0pt,text width=90mm,align=left,text=Blue,font=\fontsize{10}{12.2}\selectfont\bfseries] at (208,-64) {RESULTADO};
\node[anchor=north west,inner sep=0pt,text width=91mm,align=left,text=Blue,font=\fontsize{51}{62.2}\selectfont\bfseries] at (207,-72) {121};
\node[anchor=north west,inner sep=0pt,text width=90mm,align=left,text=Navy,font=\fontsize{14}{17.1}\selectfont\bfseries] at (208,-95) {niños por distrito};
\node[anchor=north west,inner sep=0pt,text width=90mm,align=left,text=Gray,font=\fontsize{11.5}{14.0}\selectfont] at (208,-108) {n calculado: 120,994423};
\node[anchor=north west,inner sep=0pt,text width=90mm,align=left,text=Navy,font=\fontsize{16}{19.5}\selectfont\bfseries] at (208,-120) {242 niños en total};
\node[anchor=north west,inner sep=0pt,text width=90mm,align=left,text=Gray,font=\fontsize{11}{13.4}\selectfont] at (208,-134) {Potencia con n: 80,00 \%};
\fill[Ice,rounded corners=3mm] (14,-151) rectangle (306,-165);
\node[anchor=north west,inner sep=0pt,text width=282mm,align=left,text=Navy,font=\fontsize{12}{14}\selectfont] at (19,-154.5) {pwr.2p.test devuelve n por distrito: 121 + 121 = 242 niños. No confundas n con el total.};
\end{scope}\end{tikzpicture}\mbox{}\newpage
\noindent\begin{tikzpicture}[remember picture,overlay,x=1mm,y=1mm]\begin{scope}[shift={(current page.north west)}]
\fill[Pale,rounded corners=0mm] (0,-0) rectangle (320,-180);
\node[anchor=north west,inner sep=0pt,text width=265mm,align=left,text=Blue,font=\fontsize{10}{12.2}\selectfont\bfseries] at (14,-10) {PLANTEAR  /  SJL Y VMT · VARIANTE B · POTENCIA DE 90 \%};
\node[anchor=north west,inner sep=0pt,text width=24mm,align=left,text=Navy,font=\fontsize{10}{12.2}\selectfont\bfseries] at (286,-10) {16 / 20};
\draw[Line,line width=.4pt] (14,-169) -- (306,-169);
\node[anchor=north west,inner sep=0pt,text width=250mm,align=left,text=Gray,font=\fontsize{9}{11.0}\selectfont] at (14,-172) {Jason Chambe  ·  Bioestadística  ·  TAREA 2};
\node[anchor=north west,inner sep=0pt,text width=34mm,align=left,text=Gray,font=\fontsize{9}{11.0}\selectfont\bfseries] at (274,-172) {DESLIZA >};
\node[anchor=north west,inner sep=0pt,text width=290mm,align=left,text=Navy,font=\fontsize{27}{32.9}\selectfont\bfseries] at (14,-25) {Más potencia exige más participantes};
\fill[white,rounded corners=3mm] (14,-55) rectangle (191,-147);
\fill[white,rounded corners=3mm] (197,-55) rectangle (306,-147);
\node[anchor=north west,inner sep=0pt,text width=160mm,align=left,text=Blue,font=\fontsize{10}{12.2}\selectfont\bfseries] at (21,-62) {LA COMPARACIÓN};
\node[anchor=north west,inner sep=0pt,text width=160mm,align=left,text=Navy,font=\fontsize{18}{23}\selectfont] at (21,-73) {Mantenemos las prevalencias de 43 \% y 26 \%, pero exigimos una potencia de 90 \% para la diferencia prevista.};
\node[anchor=north west,inner sep=0pt,text width=160mm,align=left,text=Blue,font=\fontsize{25}{30.5}\selectfont\bfseries] at (21,-108) {$43\,\%\quad\mathrm{vs.}\quad26\,\%$};
\node[anchor=north west,inner sep=0pt,text width=160mm,align=left,text=Navy,font=\fontsize{12.4}{15.1}\selectfont] at (21,-129) {$H_0:p_{SJL}=p_{VMT}\quad H_A:p_{SJL}\ne p_{VMT}$};
\node[anchor=north west,inner sep=0pt,text width=95mm,align=left,text=Blue,font=\fontsize{10}{12.2}\selectfont\bfseries] at (204,-62) {SUPUESTOS DEL CÁLCULO};
\node[anchor=north west,inner sep=0pt,text width=95mm,align=left,text=Navy,font=\fontsize{14}{17.1}\selectfont] at (204,-76.0) {Muestras independientes};
\node[anchor=north west,inner sep=0pt,text width=95mm,align=left,text=Navy,font=\fontsize{14}{17.1}\selectfont] at (204,-90.5) {Asignación de muestra: 1:1};
\node[anchor=north west,inner sep=0pt,text width=95mm,align=left,text=Navy,font=\fontsize{14}{17.1}\selectfont\bfseries] at (204,-105.0) {Alfa: 0,05 · Potencia: 90 \%};
\node[anchor=north west,inner sep=0pt,text width=95mm,align=left,text=Navy,font=\fontsize{14}{17.1}\selectfont] at (204,-119.5) {Alternativa: bilateral};
\fill[Ice,rounded corners=3mm] (14,-151) rectangle (306,-165);
\node[anchor=north west,inner sep=0pt,text width=282mm,align=left,text=Navy,font=\fontsize{12}{14}\selectfont] at (19,-155) {El diseño y el efecto permanecen iguales. Queremos reducir beta de 20 \% a 10 \% bajo esa alternativa.};
\node[anchor=north west,inner sep=0pt,text width=292mm,align=left,text=Gray,font=\fontsize{8.5}{10.4}\selectfont] at (14,-45) {Teoría 01, p. 47 · Teoría 02, p. 7 · Teoría 03, pp. 11 y 14};
\end{scope}\end{tikzpicture}\mbox{}\newpage
\noindent\begin{tikzpicture}[remember picture,overlay,x=1mm,y=1mm]\begin{scope}[shift={(current page.north west)}]
\fill[Pale,rounded corners=0mm] (0,-0) rectangle (320,-180);
\node[anchor=north west,inner sep=0pt,text width=265mm,align=left,text=Blue,font=\fontsize{10}{12.2}\selectfont\bfseries] at (14,-10) {CALCULAR E INTERPRETAR  /  SJL Y VMT · VARIANTE B · POTENCIA DE 90 \%};
\node[anchor=north west,inner sep=0pt,text width=24mm,align=left,text=Navy,font=\fontsize{10}{12.2}\selectfont\bfseries] at (286,-10) {17 / 20};
\draw[Line,line width=.4pt] (14,-169) -- (306,-169);
\node[anchor=north west,inner sep=0pt,text width=250mm,align=left,text=Gray,font=\fontsize{9}{11.0}\selectfont] at (14,-172) {Jason Chambe  ·  Bioestadística  ·  TAREA 2};
\node[anchor=north west,inner sep=0pt,text width=34mm,align=left,text=Gray,font=\fontsize{9}{11.0}\selectfont\bfseries] at (274,-172) {DESLIZA >};
\node[anchor=north west,inner sep=0pt,text width=292mm,align=left,text=Navy,font=\fontsize{28}{34.2}\selectfont\bfseries] at (14,-25) {Del código al tamaño que reportamos};
\node[anchor=north west,inner sep=0pt,text width=292mm,align=left,text=Blue,font=\fontsize{15}{18.3}\selectfont] at (14,-44) {$h=0{,}360193\qquad 1-\beta=0{,}90$};
\fill[Navy,rounded corners=3mm] (14,-58) rectangle (194,-147);
\fill[white,rounded corners=3mm] (200,-58) rectangle (306,-147);
\node[anchor=north west,inner sep=0pt,text width=163mm,align=left,text=Ice,font=\fontsize{10}{12.2}\selectfont\bfseries] at (21,-64) {R · PAQUETE pwr};
\node[anchor=north west,inner sep=0pt,text width=166mm,align=left,text=white,font=\fontsize{13.5}{16.9}\selectfont] at (21,-76) {\ttfamily library(pwr)\\p\_sjl\ <-\ 0.43\\p\_vmt\ <-\ 0.26\\h\ <-\ ES.h(p\_sjl,\ p\_vmt)\\calculo\ <-\ pwr.2p.test(h\ =\ h,\ sig.level\ =\ 0.05,\\\ \ power\ =\ 0.90,\ alternative\ =\ "two.sided")\\n\_por\_distrito\ <-\ ceiling(calculo\$n)\\calculo\$n\\c(por\_distrito\ =\ n\_por\_distrito,\\\ \ total\ =\ 2\ *\ n\_por\_distrito)};
\node[anchor=north west,inner sep=0pt,text width=90mm,align=left,text=Blue,font=\fontsize{10}{12.2}\selectfont\bfseries] at (208,-64) {RESULTADO};
\node[anchor=north west,inner sep=0pt,text width=91mm,align=left,text=Blue,font=\fontsize{51}{62.2}\selectfont\bfseries] at (207,-72) {162};
\node[anchor=north west,inner sep=0pt,text width=90mm,align=left,text=Navy,font=\fontsize{14}{17.1}\selectfont\bfseries] at (208,-95) {niños por distrito};
\node[anchor=north west,inner sep=0pt,text width=90mm,align=left,text=Gray,font=\fontsize{11.5}{14.0}\selectfont] at (208,-108) {n calculado: 161,977545};
\node[anchor=north west,inner sep=0pt,text width=90mm,align=left,text=Navy,font=\fontsize{16}{19.5}\selectfont\bfseries] at (208,-120) {324 niños en total};
\node[anchor=north west,inner sep=0pt,text width=90mm,align=left,text=Gray,font=\fontsize{11}{13.4}\selectfont] at (208,-134) {Potencia con n: 90,00 \%};
\fill[Ice,rounded corners=3mm] (14,-151) rectangle (306,-165);
\node[anchor=north west,inner sep=0pt,text width=282mm,align=left,text=Navy,font=\fontsize{12}{14}\selectfont] at (19,-154.5) {La muestra aumenta de 121 a 162 por distrito. Potencia no significa probabilidad de que exista un efecto.};
\end{scope}\end{tikzpicture}\mbox{}\newpage
\noindent\begin{tikzpicture}[remember picture,overlay,x=1mm,y=1mm]\begin{scope}[shift={(current page.north west)}]
\fill[Pale,rounded corners=0mm] (0,-0) rectangle (320,-180);
\node[anchor=north west,inner sep=0pt,text width=265mm,align=left,text=Blue,font=\fontsize{10}{12.2}\selectfont\bfseries] at (14,-10) {PLANTEAR  /  SJL Y VMT · VARIANTE C · MENOR DIFERENCIA};
\node[anchor=north west,inner sep=0pt,text width=24mm,align=left,text=Navy,font=\fontsize{10}{12.2}\selectfont\bfseries] at (286,-10) {18 / 20};
\draw[Line,line width=.4pt] (14,-169) -- (306,-169);
\node[anchor=north west,inner sep=0pt,text width=250mm,align=left,text=Gray,font=\fontsize{9}{11.0}\selectfont] at (14,-172) {Jason Chambe  ·  Bioestadística  ·  TAREA 2};
\node[anchor=north west,inner sep=0pt,text width=34mm,align=left,text=Gray,font=\fontsize{9}{11.0}\selectfont\bfseries] at (274,-172) {DESLIZA >};
\node[anchor=north west,inner sep=0pt,text width=290mm,align=left,text=Navy,font=\fontsize{27}{32.9}\selectfont\bfseries] at (14,-25) {Prevalencias más cercanas cuestan distinguir};
\fill[white,rounded corners=3mm] (14,-55) rectangle (191,-147);
\fill[white,rounded corners=3mm] (197,-55) rectangle (306,-147);
\node[anchor=north west,inner sep=0pt,text width=160mm,align=left,text=Blue,font=\fontsize{10}{12.2}\selectfont\bfseries] at (21,-62) {LA COMPARACIÓN};
\node[anchor=north west,inner sep=0pt,text width=160mm,align=left,text=Navy,font=\fontsize{18}{23}\selectfont] at (21,-73) {Comparamos ahora 43 \% y 33 \%, con potencia de 80 \%, alfa de 0,05 y dos muestras de igual tamaño.};
\node[anchor=north west,inner sep=0pt,text width=160mm,align=left,text=Blue,font=\fontsize{25}{30.5}\selectfont\bfseries] at (21,-108) {$43\,\%\quad\mathrm{vs.}\quad33\,\%$};
\node[anchor=north west,inner sep=0pt,text width=160mm,align=left,text=Navy,font=\fontsize{12.4}{15.1}\selectfont] at (21,-129) {$H_0:p_{SJL}=p_{VMT}\quad H_A:p_{SJL}\ne p_{VMT}$};
\node[anchor=north west,inner sep=0pt,text width=95mm,align=left,text=Blue,font=\fontsize{10}{12.2}\selectfont\bfseries] at (204,-62) {SUPUESTOS DEL CÁLCULO};
\node[anchor=north west,inner sep=0pt,text width=95mm,align=left,text=Navy,font=\fontsize{14}{17.1}\selectfont] at (204,-76.0) {Muestras independientes};
\node[anchor=north west,inner sep=0pt,text width=95mm,align=left,text=Navy,font=\fontsize{14}{17.1}\selectfont] at (204,-90.5) {Asignación de muestra: 1:1};
\node[anchor=north west,inner sep=0pt,text width=95mm,align=left,text=Navy,font=\fontsize{14}{17.1}\selectfont\bfseries] at (204,-105.0) {Alfa: 0,05 · Potencia: 80 \%};
\node[anchor=north west,inner sep=0pt,text width=95mm,align=left,text=Navy,font=\fontsize{14}{17.1}\selectfont] at (204,-119.5) {Alternativa: bilateral};
\fill[Ice,rounded corners=3mm] (14,-151) rectangle (306,-165);
\node[anchor=north west,inner sep=0pt,text width=282mm,align=left,text=Navy,font=\fontsize{12}{14}\selectfont] at (19,-155) {La diferencia se reduce de 17 a 10 puntos porcentuales. La comparación sigue siendo observacional.};
\node[anchor=north west,inner sep=0pt,text width=292mm,align=left,text=Gray,font=\fontsize{8.5}{10.4}\selectfont] at (14,-45) {Teoría 01, pp. 24 y 47 · Teoría 02, p. 7 · Teoría 03, pp. 15 y 19};
\end{scope}\end{tikzpicture}\mbox{}\newpage
\noindent\begin{tikzpicture}[remember picture,overlay,x=1mm,y=1mm]\begin{scope}[shift={(current page.north west)}]
\fill[Pale,rounded corners=0mm] (0,-0) rectangle (320,-180);
\node[anchor=north west,inner sep=0pt,text width=265mm,align=left,text=Blue,font=\fontsize{10}{12.2}\selectfont\bfseries] at (14,-10) {CALCULAR E INTERPRETAR  /  SJL Y VMT · VARIANTE C · MENOR DIFERENCIA};
\node[anchor=north west,inner sep=0pt,text width=24mm,align=left,text=Navy,font=\fontsize{10}{12.2}\selectfont\bfseries] at (286,-10) {19 / 20};
\draw[Line,line width=.4pt] (14,-169) -- (306,-169);
\node[anchor=north west,inner sep=0pt,text width=250mm,align=left,text=Gray,font=\fontsize{9}{11.0}\selectfont] at (14,-172) {Jason Chambe  ·  Bioestadística  ·  TAREA 2};
\node[anchor=north west,inner sep=0pt,text width=34mm,align=left,text=Gray,font=\fontsize{9}{11.0}\selectfont\bfseries] at (274,-172) {DESLIZA >};
\node[anchor=north west,inner sep=0pt,text width=292mm,align=left,text=Navy,font=\fontsize{28}{34.2}\selectfont\bfseries] at (14,-25) {Del código al tamaño que reportamos};
\node[anchor=north west,inner sep=0pt,text width=292mm,align=left,text=Blue,font=\fontsize{15}{18.3}\selectfont] at (14,-44) {$h=2\arcsin\sqrt{0{,}43}-2\arcsin\sqrt{0{,}33}$};
\fill[Navy,rounded corners=3mm] (14,-58) rectangle (194,-147);
\fill[white,rounded corners=3mm] (200,-58) rectangle (306,-147);
\node[anchor=north west,inner sep=0pt,text width=163mm,align=left,text=Ice,font=\fontsize{10}{12.2}\selectfont\bfseries] at (21,-64) {R · PAQUETE pwr};
\node[anchor=north west,inner sep=0pt,text width=166mm,align=left,text=white,font=\fontsize{13.5}{16.9}\selectfont] at (21,-76) {\ttfamily library(pwr)\\p\_sjl\ <-\ 0.43\\p\_vmt\ <-\ 0.33\\h\ <-\ ES.h(p\_sjl,\ p\_vmt)\\calculo\ <-\ pwr.2p.test(h\ =\ h,\ sig.level\ =\ 0.05,\\\ \ power\ =\ 0.80,\ alternative\ =\ "two.sided")\\n\_por\_distrito\ <-\ ceiling(calculo\$n)\\calculo\$n\\c(por\_distrito\ =\ n\_por\_distrito,\\\ \ total\ =\ 2\ *\ n\_por\_distrito)};
\node[anchor=north west,inner sep=0pt,text width=90mm,align=left,text=Blue,font=\fontsize{10}{12.2}\selectfont\bfseries] at (208,-64) {RESULTADO};
\node[anchor=north west,inner sep=0pt,text width=91mm,align=left,text=Blue,font=\fontsize{51}{62.2}\selectfont\bfseries] at (207,-72) {369};
\node[anchor=north west,inner sep=0pt,text width=90mm,align=left,text=Navy,font=\fontsize{14}{17.1}\selectfont\bfseries] at (208,-95) {niños por distrito};
\node[anchor=north west,inner sep=0pt,text width=90mm,align=left,text=Gray,font=\fontsize{11.5}{14.0}\selectfont] at (208,-108) {n calculado: 368,284775};
\node[anchor=north west,inner sep=0pt,text width=90mm,align=left,text=Navy,font=\fontsize{16}{19.5}\selectfont\bfseries] at (208,-120) {738 niños en total};
\node[anchor=north west,inner sep=0pt,text width=90mm,align=left,text=Gray,font=\fontsize{11}{13.4}\selectfont] at (208,-134) {Potencia con n: 80,08 \%};
\fill[Ice,rounded corners=3mm] (14,-151) rectangle (306,-165);
\node[anchor=north west,inner sep=0pt,text width=282mm,align=left,text=Navy,font=\fontsize{12}{14}\selectfont] at (19,-154.5) {El resultado es 369 por distrito, 738 en total. El tamaño de muestra no elimina la confusión entre poblaciones.};
\end{scope}\end{tikzpicture}\mbox{}\newpage
\noindent\begin{tikzpicture}[remember picture,overlay,x=1mm,y=1mm]\begin{scope}[shift={(current page.north west)}]
\fill[Navy,rounded corners=0mm] (0,-0) rectangle (320,-180);
\node[anchor=north west,inner sep=0pt,text width=265mm,align=left,text=Ice,font=\fontsize{10}{12.2}\selectfont\bfseries] at (14,-10) {GUARDA EL TUTORIAL Y PRUEBA LAS VARIANTES};
\node[anchor=north west,inner sep=0pt,text width=24mm,align=left,text=white,font=\fontsize{10}{12.2}\selectfont\bfseries] at (286,-10) {20 / 20};
\draw[Gray,line width=.4pt] (14,-169) -- (306,-169);
\node[anchor=north west,inner sep=0pt,text width=250mm,align=left,text=Ice,font=\fontsize{9}{11.0}\selectfont] at (14,-172) {Jason Chambe  ·  Bioestadística  ·  TAREA 2};
\node[anchor=north west,inner sep=0pt,text width=34mm,align=left,text=Ice,font=\fontsize{9}{11.0}\selectfont\bfseries] at (274,-172) {DESLIZA >};
\node[anchor=north west,inner sep=0pt,text width=292mm,align=left,text=white,font=\fontsize{30}{36.6}\selectfont\bfseries] at (14,-26) {La muestra depende de la pregunta};
\fill[Blue,rounded corners=3mm] (14,-55) rectangle (108,-102);
\node[anchor=north west,inner sep=0pt,text width=82mm,align=left,text=white,font=\fontsize{12}{14.6}\selectfont\bfseries] at (20,-63) {MENOR DIFERENCIA};
\node[anchor=north west,inner sep=0pt,text width=82mm,align=left,text=white,font=\fontsize{17}{20.7}\selectfont] at (20,-76) {Más muestra para distinguirla.};
\fill[Blue,rounded corners=3mm] (113,-55) rectangle (207,-102);
\node[anchor=north west,inner sep=0pt,text width=82mm,align=left,text=white,font=\fontsize{12}{14.6}\selectfont\bfseries] at (119,-63) {MAYOR POTENCIA};
\node[anchor=north west,inner sep=0pt,text width=82mm,align=left,text=white,font=\fontsize{17}{20.7}\selectfont] at (119,-76) {Más muestra para detectar el efecto.};
\fill[Blue,rounded corners=3mm] (212,-55) rectangle (306,-102);
\node[anchor=north west,inner sep=0pt,text width=82mm,align=left,text=white,font=\fontsize{12}{14.6}\selectfont\bfseries] at (218,-63) {OTRO DISEÑO};
\node[anchor=north west,inner sep=0pt,text width=82mm,align=left,text=white,font=\fontsize{17}{20.7}\selectfont] at (218,-76) {Otra comparación y otro cálculo.};
\node[anchor=north west,inner sep=0pt,text width=290mm,align=left,text=white,font=\fontsize{21}{25.6}\selectfont\bfseries] at (14,-110) {Descarga el QMD, el PDF y los scripts de R.};
\node[anchor=north west,inner sep=0pt,text width=290mm,align=left,text=Ice,font=\fontsize{15}{18.3}\selectfont\bfseries] at (14,-126) {github.com/chambejason951-collab/BIOESTADISTICA};
\node[anchor=north west,inner sep=0pt,text width=290mm,align=left,text=Ice,font=\fontsize{12}{14.6}\selectfont] at (14,-138) {Carpeta TAREA2 · 24 variantes auditadas · Resultados reproducibles};
\node[anchor=north west,inner sep=0pt,text width=292mm,align=left,text=Ice,font=\fontsize{10}{12.2}\selectfont] at (14,-151) {Basado en Teoría 01, 02 y 03 y Muestra.qmd de Antonio M. Quispe. Documentación de R y pwr. Referencias completas en el tutorial.};
\end{scope}\end{tikzpicture}\mbox{}\end{document}
